\clearpage
\SourceLeaf{241}%
\begin{center}
{\Large\textsc{[Bolyai und Lobatschewski.]}}



\emph{[\textup{\textsc{Gerling}} an \textup{\textsc{Gauss}}. Marburg, 18.--21. Dezember 1843.]}
\end{center}


\{\mbox{.\ .\ .\ .\ .\ .\ .} Ein paar ungarische Stipendiaten aus Maros Vásárhely brachten ein
Vorstellungsschreiben der dortigen Universität mit. Es fiel mir auf, weil ich unter
den übrigen Unterschriften auch den Namen \textsc{Wolfgangus Bolyai} fand. In
diesen Tagen haben sie mir auch einen Privatbrief eines anderen dortigen
Professors übergeben, der früher meine Vorlesungen besucht hatte. Ich nutzte die Gelegenheit,
mich nach W. B. zu erkundigen, um Ihnen von ihm zu berichten. Sie schildern ihn als einen
ziemlich beleibten, rüstigen und freundlichen alten Mann, der inmitten seiner
Gärten ein heiteres Alter zu verbringen scheint. Er sei ununterbrochen als Hochschullehrer
tätig und nur dadurch eingeschränkt, dass es den meisten Zuhörern zu
schwerfalle, ihm zu folgen. Sein Sohn (von dem Sie mir mitteilten, dass er ein
gutes Buch geschrieben habe) sei als erst 32-jähriger Offizier im Rang eines Rittmeisters pensioniert
worden und lebe, wenn ich richtig verstanden habe, bei seinem Vater. --- Mutmaßlich
hat er also mit mathematischen Ketzereien Anstoß erregt. --- Jener ehemalige
Hörer hatte mir damals versprochen, dafür zu sorgen, dass ich die Bücher von
Vater und Sohn hierher bekomme (denn auch vom Vater habe ich einmal ein Buch bei Ihnen gesehen, das
Sie interessant nannten). Er scheint dies aber gnädig vergessen zu haben,
denn sein Brief enthält kein Wort davon. Auch kann ich die Titel nirgends ausfindig machen
(nicht einmal in Berlin). --- Da die gerade angekommenen
jungen Leute nun anbieten, sich um die Sache zu kümmern, möchte ich Sie bitten,
mir diese Titel bei Gelegenheit aufschreiben zu lassen. \mbox{.\ .\ .\ .\ .\ .\ .}\}


